\section{Projection onto phonon modes}
\label{sec:phonons}

For a phonon branch $\nu$, let $e_{\kappa\mu,\nu}(\qq)$ denote the eigenvector in a declared normalization. The zero-point displacement is
\begin{equation}
\xi_{\kappa\mu,\nu}(\qq)
=\sqrt{\frac{\hbar}{2M_\kappa\omega_{\qq\nu}}}
 e_{\kappa\mu,\nu}(\qq).
\label{eq:zpm}
\end{equation}
The Cartesian torque kernel projects to
\begin{equation}
V_{\ell a,\nu}(\qq)
=\sum_{\kappa\mu}
K_{\ell a,\kappa\mu}(\qq)
\xi_{\kappa\mu,\nu}(\qq).
\label{eq:phonon-proj}
\end{equation}
If the electron-phonon vertex is already mode normalized,
\begin{equation}
g_\nu(\kk,\qq)
=\sum_{\kappa\mu}\xi_{\kappa\mu,\nu}(\qq)
 g_{\kappa\mu}(\kk,\qq),
\label{eq:mode-g}
\end{equation}
then \cref{eq:zpm} must not be applied again.

Near an acoustic mode at $\qq=0$, the factor $1/\sqrt{\omega}$ makes a translational-sum-rule error visible. The Cartesian response to a rigid translation must therefore be checked before phonon projection. An exactly zero or imaginary frequency is reported rather than regularized silently.
